% compendium/execution-trading-costs/futures-execution.tex
% Entries migrated on 2026-10-15 from the legacy futures-market-microstructure file (slice 19), owned by
% futures-execution; dispositions and source checks in notes/migration.md, the blocks verbatim in
% notes/artefacts/futures-execution/derivations.tex.

\section{Futures: venue conventions, off-exchange trades, margin and price discovery}

\begin{qfconvention}[CME Globex trading conventions: a centralised market with permitted exceptions, matching algorithms, the pre-open and opening match, and the E-mini S\&P 500 daily settlement window]\label{conv:cme-globex-trading}
\qffield{Statement} \emph{Centralisation.} Harris reported in 2002 that most futures markets require all trades in their contracts to be arranged within their trading systems, floor or screen. The rules prevent brokers from directing orders to confederates at prices worse than the exchange's, and futures markets had historically had no block-trading procedures or cumbersome ones \citep[chap.~``Block Traders'', sec.~15.6.2, p.~334]{alpha-harris}. CME's current rulebook keeps the principle and names the exceptions: block trades, and exchanges for related positions, which are ``permitted to be executed outside of the Exchange's centralized market in accordance with the requirements of this rule'' \citep[Rules 526 and 538]{cme-rulebook-ch5}; see \hyperref[conv:futures-block-efrp]{futures block trades and EFRPs}.

\emph{Matching.} On CME Globex a matching algorithm allocates quantity when an aggressor order meets one or more resting orders. The supported algorithms, identified by tag 1142, are Allocation (A), FIFO (F), FIFO with lead market maker (T), FIFO with top order and lead market maker (S), Pro-Rata (C), Configurable (K), Threshold Pro-Rata (O), Threshold Pro-Rata with lead market maker (Q) and an Institutional Prioritization algorithm (V). Under FIFO, price and time are the only criteria; an order loses its priority and is requeued if its quantity is increased, its price changed or its account changed. As one named example, SOFR futures (but not SOFR packs and bundles, which match FIFO) use a pro-rata allocation. A ``top'' order that betters the market has priority, then remaining quantity is allocated pro rata to resting orders at the price; orders entered in the pre-open or at the indicative opening price are matched by price and time \citep[page ``Supported Matching Algorithms'', updated 11 February 2025]{ex-cme-matching}.

\emph{Opening.} In the Pre-Open and Reserve states orders can be entered or modified, but no matching occurs, so the book can lock or cross. The trading engine computes and publishes an indicative opening price and quantity, ``a probable price and quantity at which the market will open or re-open'', for outright futures, futures spreads, options spreads and options series \citep[updated 11 February 2025]{cme-iop-wiki}.

\emph{Daily settlement, one contract.} For trade dates from 26 October 2020, the E-mini S\&P 500's lead month settles to the volume-weighted average price of its Globex trades from 14:59:30 to 15:00:00 Central Time. If it does not trade, it settles to the bid/ask midpoint over that window; the second month settles through the lead–second spread trades in the same window \citep[Exhibit~A]{cme-es-settlement-2020}.

\qffield{Conditions} The matching and opening pages are living documentation, checked on 15 October 2026; the settlement procedure is stated as filed in 2020, because the current procedure page could not be retrieved. Which algorithm applies to a given contract is set per product and was checked here only for SOFR futures. Other exchanges, and other CME contracts' settlement procedures, were not checked.

\qffield{Practice} Interpretation: the pathway's checkpoint needs a contract's algorithm before estimating passive fill probability. Under FIFO a resting order's fill depends on its queue position; under pro-rata allocation it depends on its size relative to the quantity at the price, which rewards larger posted size (\hyperref[meth:passive-order-split]{passive order split} for the FIFO queue model).

\qffield{Pitfalls} The legacy caveat \emph{No Call Auction Norm} said many futures markets do not rely on equity-style opening and closing call auctions. For CME Globex this is wrong at the open: the pre-open with an indicative opening price, followed by the opening match, is a call-auction-like uncross. For the E-mini S\&P 500 close it is right only in the narrower sense that the daily settlement is a 30-second VWAP of continuous trading, not a closing auction. Its claim that participants instead ``agree on liquid trading hours'' is not in the sources checked. The legacy principle \emph{Centralised Futures Trading} gave margining, clearing, delivery rules and fungibility as the reasons fragmentation is limited. The sources checked give exchange rules requiring on-exchange trading, with defined exceptions; the economic reasons were not checked. The legacy notes' ``highly electronic'' describes CME Globex today; Harris's 2002 account still includes floor trading.

\qffield{See also} \hyperref[conv:futures-block-efrp]{futures block trades and EFRPs}; \hyperref[mech:futures-margin-settlement]{futures margin and settlement}; \hyperref[res:spread-determinants]{spread determinants} (large ticks; owned by \texttt{execution-market-making}).

\qffield{Further reading} The legacy blocks it replaces, verbatim, are in \nolinkurl{derivations.tex} in the folder \nolinkurl{notes/artefacts/futures-execution/}.

\qffield{Index} CME Globex; matching algorithm; FIFO; pro-rata allocation; top order; lead market maker; tag 1142; SOFR futures; pre-open; indicative opening price; opening match; call auction; daily settlement price; settlement VWAP; E-mini S\&P 500; centralised market

\qffield{Status} Verification: partly checked 2026-10-15 (against the CME Client Systems Wiki pages ``Supported Matching Algorithms'' and ``Indicative Opening Price and Quantity'', both updated 11 February 2025; CME Rulebook chap.~5, Rules 526 and 538; the CFTC filing of 22 September 2020 for the E-mini S\&P 500 settlement; and the text of \texttt{alpha-harris}, 2002 Oxford printing, p.~334). Stability: draft. Currency: venue documentation as retrieved on 15 October 2026; settlement as of trade date 26 October 2020; Harris as of 2002.
\end{qfconvention}

\begin{qfconvention}[Futures block trades and exchanges for related positions at CME Group: eligibility, the related position, reporting, and the prohibition of transitory EFRPs]\label{conv:futures-block-efrp}
\qffield{Statement} \emph{Block trades} (CME Rule 526, current rulebook \citep[Rule~526]{cme-rulebook-ch5}).
\begin{itemize}
\item The exchange designates the products in which blocks are permitted and their minimum quantities, and orders may not be aggregated to reach them except by eligible advisers (526.A).
\item Each party must be an Eligible Contract Participant under the Commodity Exchange Act; registered or exempt commodity trading advisers and investment advisers with more than \$25 million under management, and similar foreign persons, may act for their customers (526.B, I, J).
\item A customer must have specified that its order be executed as a block (526.C).
\item The price must be ``fair and reasonable'' given the block's size, other trades in the contract, related markets and the circumstances (526.D).
\item Blocks do not trigger conditional orders such as stops or otherwise affect orders in the regular market (526.E).
\item The seller or broker reports the block within the period and in the manner the exchange specifies, and the exchange publishes blocks separately from regular-market trades (526.F); blocks are reported to the clearing house, and records are kept (526.G--H).
\end{itemize}
\emph{Exchanges for related positions} (EFRPs; CME Rule 538 \citep[Rule~538]{cme-rulebook-ch5}).
\begin{itemize}
\item An EFRP is a privately negotiated off-exchange execution of a futures or options contract together with the simultaneous, opposite execution of an equivalent quantity of the related cash product, by-product, related product or OTC derivative. The types are the exchange of futures for physical (EFP), for risk (EFR, against an OTC swap or other derivative) and of option for option (EOO).
\item Each party is buyer on one side and seller on the other, and each side's two legs must be for accounts with the same beneficial ownership (538.A). Opposing accounts must be independently controlled (538.B).
\item The related position must be the underlying cash commodity, a by-product, a related product or an OTC derivative with a reasonable degree of price correlation, never a futures or option on futures. It requires a bona fide transfer of ownership or a binding contract (538.C).
\item Quantities must be approximately equivalent, with hedge ratios allowed (538.E). Prices must be on the contract's increments and commercially reasonable (538.F).
\item Each EFRP is submitted to the clearing house within the specified time, with the time its terms were agreed, and is accepted only once it has cleared and the first settlement variation and performance bond are confirmed (538.I). Records are kept (538.H).
\end{itemize}
\emph{Transitory EFRPs.} Rule 538.C forbids an EFRP whose execution is contingent on another EFRP or related-position transaction between the parties, where the transactions offset the related positions without market risk that is material in context. CME Group's advisory notice in force from 7 August 2026 defines these as \emph{transitory EFRPs}. It answers ``No'' to whether they are permitted in any CME Group product: the time between the transactions is a factor, but the test is whether they ``have integrity as independent transactions exposed to market risk'' \citep[pp.~4 and~10, Q14]{cme-ra2607-5}. Immediately offsetting FX EFPs and EFRs under Rule 538.K are expressly permitted. They are not treated as transitory, because the offsetting cash or NDF transaction is not contingent on the EFRP (Q28).

\qffield{Conditions} Product thresholds and reporting deadlines are set in separate exchange notices, not checked here. The rules are those of CME Group's U.S. exchanges (CME, CBOT, NYMEX and COMEX); other venues were not checked. History, kept apart from current rules: Harris described futures block rules in 2002 as limited to eligible sophisticated traders at fair and reasonable prices. He described EFPs as originally developed for agricultural hedgers, and as used at CME in currency futures immediately after the close, where exchanges called them ``transitory EFPs'' because the floor traders bought or sold the cash currency at the same time \citep[chap.~``Block Traders'', sec.~15.6.2, p.~334]{alpha-harris}. CME Group's 2019 notice records that transitory EFRPs had been permitted in CME foreign currency, NYMEX energy and COMEX and NYMEX metals products until 4 August 2014 and have been prohibited since \citep[p.~1]{cme-ra1911-5}.

\qffield{Practice} Interpretation: an equity block cross interacts with the exchange book under precedence rules (\hyperref[def:block-trade]{block trades}, Harris's 2002 equity rules). A futures block is negotiated away from the book, does not trigger the book's conditional orders and is published separately. The two regimes should not be mixed.

\qffield{Pitfalls} The legacy taxonomy listed a ``transitory EFP: short-lived EFP used when the book is closed or thin''. The term is verified, but only as history: it matches Harris's 2002 account, and at CME Group such EFRPs have been prohibited since August 2014. The immediately offsetting FX EFRP of Rule 538.K is a different, permitted structure. Its EFP definition (``exchange cash or physical exposure for offsetting futures'') matches Rule 538 in substance; its closing sentence (``move large risk off-book, subject to eligibility and reporting rules'') fits Rules 526 and 538.

\qffield{See also} \hyperref[def:block-trade]{block trades} (owned by \texttt{execution-trading-costs}); \hyperref[conv:cme-globex-trading]{CME Globex trading conventions}; \hyperref[mech:futures-margin-settlement]{futures margin and settlement}.

\qffield{Further reading} The legacy block it replaces, verbatim, is in \nolinkurl{derivations.tex} in the folder \nolinkurl{notes/artefacts/futures-execution/}.

\qffield{Index} futures block trade; Rule 526; eligible contract participant; block minimum threshold; exchange for related position; EFRP; exchange for physical; EFP; exchange for risk; EFR; exchange of option for option; EOO; related position; transitory EFRP; immediately offsetting FX EFP; Rule 538

\qffield{Status} Verification: partly checked 2026-10-15 (against CME Rulebook chap.~5, Rules 526 and 538, as retrieved on 15 October 2026; CME Group advisory RA2607-5, effective 7 August 2026, pp.~4 and 10--17; RA1911-5 of 2019 for the history; and \texttt{alpha-harris}, 2002 Oxford printing, p.~334). Stability: draft. Currency: CME Group rules as of 15 October 2026; Harris as of 2002.
\end{qfconvention}

\begin{qfmechanism}[Futures margin and settlement at CME Clearing: performance bond, settlement variation each cycle, and the default waterfall]\label{mech:futures-margin-settlement}
\qffield{Statement} At CME Clearing, settlement for outstanding exposure is made with the clearing house in settlement cycles held ``at least once each Business Day'', at times the clearing house sets and possibly different by contract. Each clearing member, for its own and its customers' positions, pays or collects \emph{settlement variation}: the difference between this cycle's settlement price and the previous cycle's, or the trade price for a new position. Payments are final once the clearing house's bank account is debited or credited \citep[Rule~814]{cme-rulebook-ch8}. Separately, clearing members post \emph{performance bond}, initial and additional, as the clearing house sets it. It is posted gross for each omnibus customer account, in approved collateral, and retained until the positions are closed out; if a clearing member defaults on performance bond, the collateral is taken over and sold \citep[Rule~820]{cme-rulebook-ch8}. If a defaulting member's collateral does not cover its obligations, the losses are met in a stated order: CME's own contribution of \$100 million, the base guaranty fund, and then the guaranty-fund contributions of non-defaulting members \citep[Rule~802.B]{cme-rulebook-ch8}.

\qffield{Conditions} Interpretation, keeping apart the items the legacy notes merged:
\begin{itemize}
\item \emph{contract standardisation}, set in each contract's chapter (not read here);
\item \emph{initial margin}, the performance bond that covers moves between settlements and the time to close out;
\item \emph{variation settlement}, the realised gain or loss each cycle;
\item \emph{intraday calls}, which Rule 814 permits through more than one cycle a day, though the actual cycle times were not checked;
\item the \emph{FCM's} arrangements with its customers, governed by its customer agreements and by rules not read here, while Rules 814 and 820 bind clearing members;
\item \emph{residual default risk}: prices move between cycles, performance bond may fall short, and losses then pass along the waterfall of Rule 802 to the clearing house's and other members' resources.
\end{itemize}
Daily settlement bounds the exposure that can build between cycles; it does not eliminate credit risk continuously. Harris notes that futures are popular with hedgers and speculators because positions can be taken without posting large margins \citep[chap.~``Index and Portfolio Markets'', sec.~23.5, p.~492]{alpha-harris}.

\qffield{Practice} The legacy statement that the broker or FCM ``manages client margin calls, liquidation rights, and cash balances'' describes the customer side, which was not checked; the clearing-house rules above are the member side.

\qffield{Pitfalls} The legacy mechanism concluded that credit risk ``is therefore controlled continuously rather than only at final settlement''. It is controlled at each settlement cycle, at least daily, with performance bond for the moves in between and a mutualised default waterfall for what remains; ``continuously'' overstates it. The legacy taxonomy's items were checked as follows:
\begin{itemize}
\item ``fungible by expiry and specification'': contract terms were not read;
\item ``margin and default management built into the market'': Rules 802, 814 and 820;
\item ``leverage'': Harris p.~492;
\item ``hedging link to cash baskets, ETFs, options and swaps'': only partly, since Rule 538 allows cash and OTC-derivative related positions; ETFs were not checked.
\end{itemize}

\qffield{See also} \hyperref[conv:cme-globex-trading]{CME Globex trading conventions} (daily settlement prices); \hyperref[conv:futures-block-efrp]{futures block trades and EFRPs} (acceptance after first settlement variation); \hyperref[res:arbitrage-risks]{arbitrage risks} (financing of losing legs).

\qffield{Further reading} The legacy blocks it replaces, verbatim, are in \nolinkurl{derivations.tex} in the folder \nolinkurl{notes/artefacts/futures-execution/}.

\qffield{Index} futures margin; performance bond; initial margin; variation margin; settlement variation; settlement cycle; clearing house; clearing member; FCM; guaranty fund; default waterfall; mutualisation; credit risk; leverage

\qffield{Status} Verification: partly checked 2026-10-15 (against CME Rulebook chap.~8, Rules 802.B, 814 and 820, as retrieved on 15 October 2026, and \texttt{alpha-harris}, 2002 Oxford printing, p.~492). Stability: draft. Currency: CME rules as of 15 October 2026.
\end{qfmechanism}

\begin{qfevidence}[Price discovery between index futures and their cash and ETF markets: the E-mini's leadership in 2000, and what explains leadership]\label{ev:index-futures-price-discovery}
\qffield{Statement} Hasbrouck studies price discovery across U.S. equity index markets at up to one-second resolution, using vector error-correction models and information shares, for 1 March to 31 May 2000. He finds:
\begin{itemize}
\item for the S\&P 500 and Nasdaq-100, whose index markets comprised an ETF, a floor-traded futures contract and an electronic E-mini, most price discovery occurred in the E-mini, whose prices generally led the regular futures and the ETFs; the ETFs' contributions were statistically significant but small;
\item for the S\&P 400 MidCap, with only an ETF and regular futures, both contributed substantially and the ETF appeared to dominate;
\item the S\&P 500 sector ETFs contributed little to price discovery in the index.
\end{itemize}
\citep[NYU working paper FIN-00-046, draft of 16 November 2000, abstract and pp.~4--5, 13 and~34]{ex-hasbrouck-2003}. The earlier literature he reviews found that the dominant influence ran from futures to cash, with weaker but measurable effects in the reverse direction (p.~5).

\qffield{Conditions} The evidence is for three U.S. indexes in three months of 2000, when the regular contracts traded on the CME floor during regular hours and on Globex otherwise, and the E-minis on Globex (p.~6). It is from a preliminary working-paper draft. The catalogue describes the published 2003 version's S\&P 400 result as shared between the regular futures and the ETF; the published version was not read, so the difference is recorded, not resolved. The theoretical treatment of futures--cash price discovery the pathway assigns, Garbade and Silber (1983, \texttt{w5-garbade-silber}), was not available and is unread.

\qffield{Practice} Harris's reason for index-market leadership is that index traders price only index risk, while traders of the components must price every risk of their securities, so index markets discover the price of index risk faster (\hyperref[def:portfolio-trading]{portfolio trading}, p.~489).

\qffield{Pitfalls} The legacy mechanism said futures lead ``when they are cheaper to trade, more leveraged, or open during information arrival'', and cash markets lead ``when constituent-level information dominates''. Harris's argument supports only the general direction for index risk. The specific conditions are plausible hypotheses not tested in the passages read, and the working paper does not attribute the E-mini's leadership to cost or leverage. Leadership in information shares is a statistical measure of where innovations first appear, not a causal claim about who is informed.

\qffield{See also} \hyperref[def:portfolio-trading]{portfolio trading}, on why index markets discover the price of index risk first, owned by the section overview \texttt{execution-trading-costs}; \hyperref[mech:etf-creation-redemption]{ETF creation and redemption}; \hyperref[def:arbitrage]{arbitrage} (the basis and delivery arbitrage).

\qffield{Further reading} The legacy block it replaces, verbatim, is in \nolinkurl{derivations.tex} in the folder \nolinkurl{notes/artefacts/futures-execution/}. Garbade and Silber (1983) for the theory.

\qffield{Index} price discovery; lead--lag; index futures; E-mini; ETF; information share; vector error correction model; cointegration; basis; futures leads cash; S\&P 500; Nasdaq-100; S\&P 400

\qffield{Status} Verification: partly checked 2026-10-15 (against \texttt{ex-hasbrouck-2003} in its NYU working-paper draft of 16 November 2000, abstract and pp.~4--5, 13 and~34, text layer, not the Journal of Finance version; Garbade and Silber unread). Stability: draft. Currency: evidence for March--May 2000.
\end{qfevidence}
